\documentclass[11pt]{article}
\usepackage{mathblatt}
\newcommand{\bereich}[1]{\par\addvspace{14pt}\noindent{\large\bfseries #1}\par\nobreak\vspace{4pt}}
\newcommand{\bl}[1]{\par\noindent #1\par\vspace{3pt}}
\newcommand{\blrand}[1]{\par\noindent{\color{mbgrau}#1}\par\vspace{3pt}}
\newcommand{\blhier}[1]{\par\noindent\llap{$\blacktriangleright$\hspace{0.5em}}\textbf{#1}\par\vspace{3pt}}
\newcommand{\stufe}[1]{\par\noindent\hspace*{1.5em}{\small #1}\par\vspace{2pt}}
\newcommand{\pkt}{\textperiodcentered{}}
\begin{document}

\pfheftstil
\pflaufkopf[Abitur GK \pkt{} Prüfung]{Kurvenuntersuchung}{GK-KU10}\pfseitevon
\pfheftkopf{Kurvenuntersuchung}{}
\blhier{Nullstellen und Achsenschnittpunkte berechnen}
\stufe{Funktionswert und Nullstelle durch Einsetzen}
\stufe{Nullstellen ganzrationaler Funktionen}
\stufe{Achsenschnittpunkte bei e-Funktionen}
\bl{Extrempunkte berechnen}
\bl{Wendepunkte berechnen}
\bl{Monotonie und Krümmung aus der Ableitung begründen}
\bl{Grenzverhalten angeben}
\bl{Symmetrie am Term begründen}
\blrand{Graph verschieben, spiegeln, strecken}
\blrand{Maße im Sachzusammenhang aus dem Graphen berechnen}
\blrand{Gleichungen und Ungleichungen zwischen Funktionen lösen}
\blrand{Graph skizzieren und Verlauf beschreiben}
\blrand{Funktionsgleichung aus Bedingungen aufstellen}
\blrand{Extremalproblem: Zielfunktion aufstellen und maximieren}

\par\addvspace{14pt}\noindent Weiter: Ableitung + Tangente\par
\end{document}